\documentclass[10pt,a4paper]{article}

\usepackage[utf8]{inputenc}
\usepackage[T1]{fontenc}
\usepackage[a4paper,margin=0.6in,top=0.52in,bottom=0.52in]{geometry}
\usepackage{amsmath}
\usepackage{tikz}
\usetikzlibrary{shapes,arrows,positioning,calc,fit,backgrounds}
\usepackage{xcolor}
\usepackage{lmodern}
\usepackage{booktabs}
\usepackage{tabularx}
\usepackage{enumitem}
\usepackage{titlesec}
\usepackage{fancyhdr}
\usepackage{listings}
\usepackage{microtype}
\usepackage[colorlinks=true,linkcolor=primary,urlcolor=secondary]{hyperref}

% =========================================================
% Palette & Styling (Matching lab_9.tex)
% =========================================================
\definecolor{primary}{RGB}{24, 43, 73}        % Deep clean navy
\definecolor{secondary}{RGB}{14, 116, 144}   % Teal accent
\definecolor{darkslate}{RGB}{30, 41, 59}     % Normal text color
\definecolor{boxbg}{RGB}{248, 250, 252}      % Soft light box
\definecolor{bordergray}{RGB}{226, 232, 240} % Light border

% Terminal Color Scheme: Modern Light Grey Theme
\definecolor{termbg}{RGB}{246, 248, 250}     % Light grey terminal body
\definecolor{termbar}{RGB}{229, 231, 235}    % Medium light grey title bar
\definecolor{termborder}{RGB}{209, 213, 219} % Crisp border line
\definecolor{termtext}{RGB}{31, 41, 55}      % Dark charcoal text
\definecolor{termuser}{RGB}{22, 101, 52}     % Dark green for user
\definecolor{termhost}{RGB}{21, 128, 61}     % Green for host
\definecolor{termpath}{RGB}{29, 78, 216}     % Blue for directory path
\definecolor{termcmd}{RGB}{17, 24, 39}       % Bold black for commands
\definecolor{termout}{RGB}{55, 65, 81}       % Clear dark grey for outputs
\definecolor{termdir}{RGB}{37, 99, 235}      % Blue for directories
\definecolor{termfile}{RGB}{31, 41, 55}      % Dark for regular files
\definecolor{termhidden}{RGB}{107, 114, 128} % Muted grey
\definecolor{termsize}{RGB}{180, 83, 9}      % Amber for sizes
\definecolor{termperm}{RGB}{79, 70, 229}     % Indigo for permissions

% Code Listing Colors
\definecolor{codebg}{RGB}{248, 250, 252}
\definecolor{codeborder}{RGB}{218, 224, 233}
\definecolor{codekw}{RGB}{14, 116, 144}
\definecolor{codestr}{RGB}{22, 101, 52}
\definecolor{codecomment}{RGB}{100, 116, 139}
\definecolor{codenum}{RGB}{148, 163, 184}

% Table helpers
\newcolumntype{Y}{>{\centering\arraybackslash}X}

% =========================================================
% Section Headings
% =========================================================
\titleformat{\section}
  {\color{primary}\large\bfseries\sffamily\raggedright}
  {\thesection.}{0.5em}{}
\titlespacing*{\section}{0pt}{7pt}{2pt}

\titleformat{\subsection}
  {\color{secondary}\normalsize\bfseries\sffamily\raggedright}
  {\thesubsection}{0.5em}{}
\titlespacing*{\subsection}{0pt}{5pt}{2pt}

% =========================================================
% Headers and Footers
% =========================================================
\pagestyle{fancy}
\fancyhf{}
\renewcommand{\headrulewidth}{0.4pt}
\renewcommand{\footrulewidth}{0pt}
\fancyhead[L]{\small\sffamily\color{darkslate} Operating Systems Lab}
\fancyhead[R]{\small\sffamily\color{darkslate} Lab Report: FCFS Scheduling Algorithm}
\fancyfoot{}

\fancypagestyle{firstpage}{%
  \fancyhf{}
  \renewcommand{\headrulewidth}{0pt}
  \renewcommand{\footrulewidth}{0pt}
}

% =========================================================
% Terminal Window Macro
% =========================================================
\newcommand{\termprompt}[1][\textasciitilde]{%
  \textbf{\textcolor{termuser}{ratul}}\textcolor{termtext}{@}\textbf{\textcolor{termhost}{firefly}}\textcolor{termtext}{:}\textbf{\textcolor{termpath}{#1}}\textcolor{termcmd}{\$}\ %
}

\newcommand{\terminalwindow}[3][\dimexpr\linewidth-0.6pt\relax]{%
\par\vspace{2pt}%
{\centering
\begin{tikzpicture}
  % 1. Terminal text on main layer
  \node[
      text=termtext,
      font=\ttfamily\fontsize{7.2}{9.2}\selectfont,
      align=left,
      text width=\dimexpr#1-1.02cm\relax,
      inner xsep=6pt,
      inner ysep=6pt,
      anchor=north west
  ] (body) at (0.32cm, -0.56cm) {#3};

  % Coordinates
  \coordinate (topleft) at (0, 0);
  \coordinate (bottomright) at ($(body.south -| #1, 0) + (0, -5pt)$);
  \coordinate (titlecenter) at ($0.5*(topleft) + 0.5*(#1, 0) + (0, -0.28cm)$);

  % 2. Background and Title bar on background layer
  \begin{scope}[on background layer]
    \fill[termbg, rounded corners=6pt] (topleft) rectangle (bottomright);

    \begin{scope}
      \clip[rounded corners=6pt] (topleft) rectangle (bottomright);
      \fill[termbar] (topleft) rectangle (#1, -0.56cm);
      \draw[termborder, line width=0.5pt] (0, -0.56cm) -- (#1, -0.56cm);
    \end{scope}

    \draw[termborder, line width=0.6pt, rounded corners=6pt] (topleft) rectangle (bottomright);
  \end{scope}

  % 3. Window control dots
  \fill[red!80!black] (0.36cm, -0.28cm) circle (0.07cm);
  \fill[yellow!80!black] (0.63cm, -0.28cm) circle (0.07cm);
  \fill[green!75!black] (0.90cm, -0.28cm) circle (0.07cm);

  % 4. Centered terminal title
  \node[
      text=gray!40!black,
      font=\sffamily\bfseries\fontsize{7.2}{8}\selectfont,
      anchor=center
  ] at (titlecenter) {#2};
\end{tikzpicture}\par}%
\vspace{2pt}%
}

\newcommand{\cmd}[1]{\texttt{\textbf{#1}}}

% =========================================================
% Code Listing Configuration
% =========================================================
\lstset{
  language=C++,
  backgroundcolor=\color{codebg},
  basicstyle=\ttfamily\fontsize{7.2}{9.0}\selectfont\color{darkslate},
  keywordstyle=\bfseries\color{codekw},
  commentstyle=\itshape\color{codecomment},
  stringstyle=\color{codestr},
  numbers=left,
  numberstyle=\tiny\color{codenum},
  stepnumber=1,
  numbersep=8pt,
  showstringspaces=false,
  breaklines=true,
  frame=single,
  rulecolor=\color{codeborder},
  tabsize=4,
  xleftmargin=14pt,
  framexleftmargin=14pt
}

\begin{document}
\raggedbottom
\thispagestyle{firstpage}

% =========================================================
% Title Block
% =========================================================
\begin{center}
  {\Large\bfseries\sffamily\color{primary} Experiment Title: Implementation of First Come, First Served CPU Scheduling Algorithm in C++}\\[3pt]
  {\normalsize\sffamily\textbf{Topic:} Non-Preemptive FCFS CPU Scheduling Simulation with Arrival Time in C++}
\end{center}
\vspace{-4pt}

% =========================================================
% 1. Problem
% =========================================================
\section{Problem}
Implement the non-preemptive First Come, First Served (FCFS) CPU scheduling algorithm in C++. Given the number of processes, arrival times ($AT$), and execution/burst times ($BT$), schedule the processes strictly by arrival order, handle CPU idle time, and compute Completion Time ($CT$), Turnaround Time ($TAT$), Waiting Time ($WT$), and average metrics along with a Gantt chart.

% =========================================================
% 2. Theoretical Solution
% =========================================================
\section{Theoretical Solution}
FCFS allocates the CPU to processes strictly in their order of arrival. Execution is non-preemptive: once a process gains the CPU, it runs to completion. If the CPU is idle when a process arrives, its start time equals its arrival time ($ST_i = \max(currentTime, AT_i)$).

\subsection{Formulas}
For each process $i$:
\begin{itemize}[noitemsep,topsep=1pt,leftmargin=1.2em]
  \item $CT_i = ST_i + BT_i$
  \item $TAT_i = CT_i - AT_i$
  \item $WT_i = TAT_i - BT_i$
  \item $Avg\_TAT = \frac{1}{n} \sum TAT_i, \quad Avg\_WT = \frac{1}{n} \sum WT_i$
\end{itemize}

\subsection{Algorithmic Steps}
\begin{enumerate}[noitemsep,topsep=1pt,leftmargin=1.4em]
  \item Read process count $n$, arrival times $AT$, and burst times $BT$.
  \item Sort processes by arrival time using stable sort.
  \item For each process, update $currentTime = \max(currentTime, AT_i)$, then compute $CT_i$, $TAT_i$, and $WT_i$.
  \item Display the schedule table, average metrics, and Gantt chart sequence.
\end{enumerate}

% =========================================================
% 3. Expected Input/Output (I/O) Set
% =========================================================
\section{Expected Input/Output (I/O) Set}
Expected input and scheduling results for $n = 4$ processes:

\vspace{-2pt}
\begin{table}[h!]
\centering\footnotesize
\renewcommand{\arraystretch}{1.08}
\begin{tabularx}{\linewidth}{YYYYYYY}
\toprule
\textbf{Process} & \textbf{Arrival ($AT$)} & \textbf{Execution ($BT$)} & \textbf{Start ($ST$)} & \textbf{Completion ($CT$)} & \textbf{Turnaround ($TAT$)} & \textbf{Waiting ($WT$)} \\
\midrule
\textbf{P1} & 0 & 4 & 0 & 4  & 4 & 0 \\
\textbf{P2} & 1 & 3 & 4 & 7  & 6 & 3 \\
\textbf{P3} & 2 & 1 & 7 & 8  & 6 & 5 \\
\textbf{P4} & 3 & 2 & 8 & 10 & 7 & 5 \\
\bottomrule
\end{tabularx}
\end{table}
\vspace{-6pt}

\noindent\textbf{Expected Averages:} Average Waiting Time = $3.25$, Average Turnaround Time = $5.75$.

% =========================================================
% 4. Source Code
% =========================================================
\section{Source Code}
Complete C++ implementation of the First Come, First Served scheduling algorithm:

\begin{lstlisting}
#include <bits/stdc++.h>
using namespace std;

struct Process {
    int id;
    int arrival;
    int execution;
    int start;
    int completion;
    int turnaround;
    int waiting;
};

int main() {
    int n;

    cout << "Enter the number of processes: ";
    cin >> n;

    if (cin.fail() || n <= 0) {
        cout << "Invalid number of processes!\n";
        return 1;
    }

    vector<Process> processes(n);

    // Input process details
    for (int i = 0; i < n; i++) {
        processes[i].id = i + 1;

        cout << "Enter arrival and execution time for process "
             << processes[i].id << ": ";

        cin >> processes[i].arrival >> processes[i].execution;

        if (cin.fail() ||
            processes[i].arrival < 0 ||
            processes[i].execution <= 0) {
            cout << "Invalid input! Arrival time must be "
                 << "non-negative and execution time must "
                 << "be positive.\n";
            return 1;
        }
    }

    // Sort by arrival time; preserve input order for ties
    stable_sort(processes.begin(), processes.end(),
        [](const Process& a, const Process& b) {
            return a.arrival < b.arrival;
        });

    int currentTime = 0;
    double totalWaiting = 0;
    double totalTurnaround = 0;

    // Calculate scheduling times
    for (Process& p : processes) {
        // Handle CPU idle time
        currentTime = max(currentTime, p.arrival);

        p.start = currentTime;
        p.waiting = p.start - p.arrival;

        currentTime += p.execution;

        p.completion = currentTime;
        p.turnaround = p.completion - p.arrival;

        totalWaiting += p.waiting;
        totalTurnaround += p.turnaround;
    }

    // Display results in original process order
    sort(processes.begin(), processes.end(),
        [](const Process& a, const Process& b) {
            return a.id < b.id;
        });

    cout << "\nProcess\tArrival\tExecution\tStart\tCompletion"
         << "\tTurnaround\tWaiting\n";

    for (const Process& p : processes) {
        cout << "P" << p.id << "\t"
             << p.arrival << "\t"
             << p.execution << "\t\t"
             << p.start << "\t"
             << p.completion << "\t"
             << p.turnaround << "\t\t"
             << p.waiting << '\n';
    }

    // Display average times
    cout << fixed << setprecision(2);

    cout << "\nAverage Waiting Time: "
         << totalWaiting / n << '\n';

    cout << "Average Turnaround Time: "
         << totalTurnaround / n << '\n';

    // Display Gantt chart in FCFS execution order
    stable_sort(processes.begin(), processes.end(),
        [](const Process& a, const Process& b) {
            return a.arrival < b.arrival;
        });

    cout << "\nGantt Chart:\n";

    for (const Process& p : processes) {
        if (p.start > 0 && p.start > p.arrival) {
            // The CPU may have been idle before this process.
            // Idle intervals are displayed below using timestamps.
        }

        cout << "| ";
        if (p.start > 0) {
            cout << "P" << p.id << " ";
        } else {
            cout << "P" << p.id << " ";
        }
    }
    cout << "|\n";

    // Print process start and completion timestamps
    if (!processes.empty()) {
        cout << processes[0].start;

        for (const Process& p : processes) {
            cout << "\t" << p.completion;
        }

        cout << '\n';
    }

    return 0;
}
\end{lstlisting}

% =========================================================
% 5. Execution of Code
% =========================================================
\section{Execution of Code}
The program was compiled with \cmd{g++} and executed in Linux terminal:

\terminalwindow{ratul@firefly: \textasciitilde/OS-Lab\ \textemdash\ bash}{
% Command 1: compile
\termprompt[\textasciitilde/OS-Lab]\textbf{\textcolor{termcmd}{g++ fcfs.cpp -o fcfs}}

\vspace{0.06cm}
% Command 2: run
\termprompt[\textasciitilde/OS-Lab]\textbf{\textcolor{termcmd}{./fcfs}}

\textcolor{termout}{Enter the number of processes: 4}

\textcolor{termout}{Enter arrival and execution time for process 1: 0 4}

\textcolor{termout}{Enter arrival and execution time for process 2: 1 3}

\textcolor{termout}{Enter arrival and execution time for process 3: 2 1}

\textcolor{termout}{Enter arrival and execution time for process 4: 3 2}

\vspace{0.06cm}
\textcolor{termperm}{\textbf{Process\ \ Arrival\ \ Execution\ \ Start\ \ Completion\ \ Turnaround\ \ Waiting}}

\textcolor{termout}{P1\ \ \ \ \ \ \ 0\ \ \ \ \ \ \ \ 4\ \ \ \ \ \ \ \ \ \ 0\ \ \ \ \ \ 4\ \ \ \ \ \ \ \ \ \ \ 4\ \ \ \ \ \ \ \ \ \ \ 0}

\textcolor{termout}{P2\ \ \ \ \ \ \ 1\ \ \ \ \ \ \ \ 3\ \ \ \ \ \ \ \ \ \ 4\ \ \ \ \ \ 7\ \ \ \ \ \ \ \ \ \ \ 6\ \ \ \ \ \ \ \ \ \ \ 3}

\textcolor{termout}{P3\ \ \ \ \ \ \ 2\ \ \ \ \ \ \ \ 1\ \ \ \ \ \ \ \ \ \ 7\ \ \ \ \ \ 8\ \ \ \ \ \ \ \ \ \ \ 6\ \ \ \ \ \ \ \ \ \ \ 5}

\textcolor{termout}{P4\ \ \ \ \ \ \ 3\ \ \ \ \ \ \ \ 2\ \ \ \ \ \ \ \ \ \ 8\ \ \ \ \ \ 10\ \ \ \ \ \ \ \ \ \ 7\ \ \ \ \ \ \ \ \ \ \ 5}

\vspace{0.06cm}
\textcolor{termuser}{\textbf{Average Waiting Time: 3.25}}

\textcolor{termuser}{\textbf{Average Turnaround Time: 5.75}}

\vspace{0.06cm}
\textcolor{termsize}{\textbf{Gantt Chart:}}

\textcolor{termcmd}{| P1 | P2 | P3 | P4 |}

\textcolor{termout}{0\ \ \ \ \ \ \ 4\ \ \ \ \ \ \ 7\ \ \ \ \ \ \ 8\ \ \ \ \ \ \ 10}

\vspace{0.08cm}
% Final prompt
\termprompt[\textasciitilde/OS-Lab]\textbf{\textcolor{termcmd}{\rule{0pt}{0.9em}\textunderscore}}
}

\vspace{4pt}\noindent\textbf{Verification:}\ The execution output matches the expected completion times, turnaround times, waiting times, and Gantt chart sequence, confirming correct non-preemptive FCFS scheduling.

\end{document}
